\documentclass[12pt,a4paper]{article}

% ── Packages ────────────────────────────────────────────────────────────────
\usepackage[a4paper, margin=2.5cm]{geometry}
\usepackage{graphicx}
\usepackage{hyperref}
\usepackage{hypcap}
\usepackage{bookmark}
\usepackage{amsmath,amssymb,amsfonts}
\usepackage{titlesec}
\usepackage{caption}
\usepackage{subcaption}
\usepackage{enumitem}
\usepackage{lmodern}
\usepackage{cmap}
\usepackage{xcolor}
\usepackage{float}
\usepackage{colortbl}

% Soft pastel colors for headings
\definecolor{softblue}{HTML}{6B9AC4}
\definecolor{softteal}{HTML}{7FB3B3}
\definecolor{softpurple}{HTML}{B39EB5}
\definecolor{low}{HTML}{F8E6E6}
\definecolor{mid}{HTML}{E6F2F2}
\definecolor{high}{HTML}{D6EAF8}
\definecolor{best}{HTML}{A9CCE3}

% Table and long table support
\usepackage{tabularx}
\usepackage{longtable}
\usepackage{booktabs}
\usepackage{multirow}
\usepackage{array}
\usepackage{fancyhdr}
\usepackage[T1]{fontenc}
\usepackage[utf8]{inputenc}
\usepackage{comment}
\usepackage{chngcntr}
\usepackage{placeins}
\counterwithin{figure}{section}
\counterwithin{table}{section}

% ── Hyperref setup ───────────────────────────────────────────────────────────
\hypersetup{
    colorlinks=true,
    linkcolor=black,
    urlcolor=black,
    citecolor=black
}

% ── Section formatting ───────────────────────────────────────────────────────
\titleformat{\section}{\large\bfseries\color{softblue}}{\thesection}{1em}{}
\titleformat{\subsection}{\normalsize\bfseries}{\thesubsection}{1em}{}
\titleformat{\subsubsection}{\normalsize\bfseries}{\thesubsubsection}{1em}{}

% ── Header / Footer ──────────────────────────────────────────────────────────
\pagestyle{fancy}
\fancyhf{}
\rhead{\small Neural Networks -- Assignment 3}
\lhead{\small Cairo University}
\cfoot{\thepage}
\setlength{\headheight}{14pt}
\renewcommand{\headrulewidth}{0.4pt}

% ── Column types ─────────────────────────────────────────────────────────────
\newcolumntype{C}[1]{>{\centering\arraybackslash}p{#1}}
\newcolumntype{L}[1]{>{\raggedright\arraybackslash}p{#1}}

\graphicspath{{./}}

% ══════════════════════════════════════════════════════════════════════════════
\begin{document}

% ── Title Page ───────────────────────────────────────────────────────────────
\begin{titlepage}
\centering
\vspace*{0.5cm}

{\Large\bfseries Cairo University\\}
\vspace{0.2cm}
{\large Faculty of Engineering\\}
\vspace{0.2cm}
{\normalsize Electronics and Electrical Communications Engineering\\}

\vspace{1.2cm}
\rule{0.8\textwidth}{1.5pt}\\[0.4cm]
{\LARGE\bfseries Assignment 3}\\[0.3cm]
{\large\bfseries ELC 4028 -- Artificial Neural Networks}\\[0.2cm]
{\large\bfseries Problem 3: Impact of Attention Mechanisms}\\[0.2cm]
\rule{0.8\textwidth}{1.5pt}

\vspace{1.2cm}

\begin{tabular}{|c|c|l|}
\hline
\textbf{ID} & \textbf{Section} & \textbf{Name} \\
\hline
9220399 & 2 & Shahd Gamal Mahmoud Hussein \\
\hline
9220411 & 2 & Salah El-Din Hassan Hassan Mohamed \\
\hline
9220984 & 4 & Yousef Khaled Omar Mahmoud \\
\hline
\end{tabular}

\vspace{1.2cm}

\textit{Supervised by:}\\[0.3cm]
\textit{Dr.\ Mohsen Rashwan}\\
\textit{Dr.\ Mohamed Motaz}

\vfill
\end{titlepage}

% ── Table of Contents ────────────────────────────────────────────────────────
\setcounter{tocdepth}{3}
\tableofcontents
\newpage
\listoffigures
\newpage
\listoftables
\newpage

% ══════════════════════════════════════════════════════════════════════════════
%  PART (a)
% ══════════════════════════════════════════════════════════════════════════════
\section{Part (a): CNN vs.\ Attention CNN on ReducedMNIST}

\subsection{Introduction}

This part builds two CNN classifiers on the ReducedMNIST dataset — a subset of
the standard MNIST handwritten-digit benchmark — and measures the impact of
inserting a spatial attention mechanism. The two models are:

\begin{enumerate}[leftmargin=1.5em]
    \item \textbf{Baseline CNN} — a compact two-block convolutional network with no
          attention, similar in spirit to LeNet-5.
    \item \textbf{Attention CNN} — the same architecture augmented with a
          \texttt{SpatialAttention} layer after each convolutional block.
\end{enumerate}

\subsection{Dataset}

ReducedMNIST is a down-sampled version of MNIST consisting of 28$\times$28 greyscale
images of handwritten digits (0--9). The dataset is split into training and validation
sets, with inputs normalised to $[0,1]$.

\subsection{Network Architectures}

\subsubsection{Baseline CNN}

The baseline model is a compact sequential network. Table~\ref{tab:a_baseline_arch}
shows its layer structure extracted from the model summary.

\begin{table}[H]
\centering
\caption{Baseline CNN architecture for ReducedMNIST}
\label{tab:a_baseline_arch}
\begin{tabular}{L{5cm} C{3.5cm} C{2.5cm}}
\toprule
\rowcolor{best}
\textbf{Layer (type)} & \textbf{Output Shape} & \textbf{Param \#} \\
\midrule
Conv2D 3$\times$3, 4 filters, ReLU  & (None, 26, 26, 4)  & 40 \\
\rowcolor{high}
MaxPooling2D 2$\times$2              & (None, 13, 13, 4)  & 0 \\
Conv2D 3$\times$3, 8 filters, ReLU  & (None, 11, 11, 8)  & 296 \\
\rowcolor{high}
MaxPooling2D 2$\times$2              & (None, 5, 5, 8)    & 0 \\
Flatten                              & (None, 200)         & 0 \\
\rowcolor{high}
Dense (60, ReLU)                     & (None, 60)          & 12,060 \\
Dense (10, Softmax)                  & (None, 10)          & 610 \\
\midrule
\rowcolor{best}
\multicolumn{2}{l}{\textbf{Total trainable parameters}} & \textbf{13,006} \\
\multicolumn{2}{l}{Non-trainable parameters} & 0 \\
\bottomrule
\end{tabular}
\end{table}

\subsubsection{Spatial Attention Mechanism}

The \texttt{SpatialAttention} layer generates a 2D attention mask from the feature
map by applying average-pooling and max-pooling along the channel dimension,
concatenating the results, and passing them through a $3\times3$ (or $7\times7$)
sigmoid convolution:

\begin{equation}
    \mathbf{M} = \sigma\!\left(
        \text{Conv}\!\left(
            \left[\,\text{AvgPool}_{C}(\mathbf{F});\;
                  \text{MaxPool}_{C}(\mathbf{F})\,\right]
        \right)
    \right), \qquad
    \mathbf{F}' = \mathbf{F} \odot \mathbf{M}
\end{equation}

Each \texttt{SpatialAttention} module adds \textbf{19 parameters}:
$(3\times3\times2)+1 = 19$.

\subsubsection{Attention CNN}

The attention model inserts a \texttt{SpatialAttention} layer immediately after
each convolutional layer, before pooling. Table~\ref{tab:a_attention_arch} shows
the full structure.

\begin{table}[H]
\centering
\caption{Attention CNN architecture for ReducedMNIST}
\label{tab:a_attention_arch}
\begin{tabular}{L{5.5cm} C{3.5cm} C{2cm}}
\toprule
\rowcolor{best}
\textbf{Layer (type)} & \textbf{Output Shape} & \textbf{Param \#} \\
\midrule
Conv2D 3$\times$3, 4 filters, ReLU        & (None, 26, 26, 4)  & 40 \\
\rowcolor{mid}
\textbf{SpatialAttention}                  & (None, 26, 26, 4)  & \textbf{19} \\
MaxPooling2D 2$\times$2                    & (None, 13, 13, 4)  & 0 \\
\rowcolor{high}
Conv2D 3$\times$3, 8 filters, ReLU        & (None, 11, 11, 8)  & 296 \\
\rowcolor{mid}
\textbf{SpatialAttention}                  & (None, 11, 11, 8)  & \textbf{19} \\
MaxPooling2D 2$\times$2                    & (None, 5, 5, 8)    & 0 \\
\rowcolor{high}
Flatten                                    & (None, 200)         & 0 \\
Dense (60, ReLU)                           & (None, 60)          & 12,060 \\
\rowcolor{high}
Dense (10, Softmax)                        & (None, 10)          & 610 \\
\midrule
\rowcolor{best}
\multicolumn{2}{l}{\textbf{Total trainable parameters}} & \textbf{13,044} \\
\multicolumn{2}{l}{Non-trainable parameters} & 0 \\
\bottomrule
\end{tabular}
\end{table}

\subsubsection{Parameter Overhead of Attention}

Two \texttt{SpatialAttention} modules add $2 \times 19 = \mathbf{+38}$ parameters
on top of 13,006 — an overhead of $\approx 0.29\%$.

\subsection{Training Configuration}

Both models are trained under identical conditions summarised in
Table~\ref{tab:a_hyperparams}.

\begin{table}[H]
\centering
\caption{Training hyperparameters for Part (a)}
\label{tab:a_hyperparams}
\begin{tabular}{L{5.5cm} C{5cm}}
\toprule
\rowcolor{best}
\textbf{Hyperparameter} & \textbf{Value} \\
\midrule
Optimizer         & Adam \\
\rowcolor{high}
Loss function     & Sparse categorical cross-entropy \\
Epochs            & 10 \\
\rowcolor{high}
Input image size  & 28 $\times$ 28 $\times$ 1 \\
Number of classes & 10 (digits 0--9) \\
\bottomrule
\end{tabular}
\end{table}

\subsubsection{Training Logs}

\begin{figure}[H]
    \centering
    \includegraphics[width=\linewidth]{training_baseline_a.jpg}
    \caption{Baseline CNN training log (ReducedMNIST).}
    \label{fig:a_log_baseline}
\end{figure}

\begin{figure}[H]
    \centering
    \includegraphics[width=\linewidth]{training_attention_a.jpg}
    \caption{Attention CNN training log (ReducedMNIST).}
    \label{fig:a_log_attention}
\end{figure}

\subsection{Results}

\subsubsection{Training and Validation Curves}

\begin{figure}[H]
    \centering
    \includegraphics[width=\linewidth]{baseline_curves.png}
    \caption{Baseline CNN accuracy and loss curves (ReducedMNIST).}
    \label{fig:a_curves_baseline}
\end{figure}

\begin{figure}[H]
    \centering
    \includegraphics[width=\linewidth]{attention_curves.png}
    \caption{Attention CNN accuracy and loss curves (ReducedMNIST).}
    \label{fig:a_curves_attention}
\end{figure}

\subsubsection{Quantitative Comparison}

\begin{table}[H]
\centering
\caption{Final-epoch performance: Baseline vs.\ Attention CNN (ReducedMNIST)}
\label{tab:a_results}
\begin{tabular}{L{4.5cm} C{2.8cm} C{2.8cm} C{2.3cm}}
\toprule
\rowcolor{best}
\textbf{Metric} & \textbf{Baseline CNN} & \textbf{Attention CNN} & \textbf{Difference} \\
\midrule
Train Accuracy (\%)    & 97.46 & 97.52 & \textcolor{softblue}{+0.06} \\
\rowcolor{high}
Test Accuracy (\%)     & 96.65 & 96.58 & $-$0.07 \\
Train Loss             & 0.0832 & 0.0796 & \textcolor{softblue}{$-$0.0036} \\
\rowcolor{high}
Test Loss              & 0.1032 & 0.1166 & $+$0.0134 \\
Training Time (s)      & 11.75 & 21.65 & $+$9.90 \\
\rowcolor{high}
Train-Test Acc.\ Gap   & 0.81\% & 0.94\% & $+$0.13 pp \\
Model Parameters       & 13,006 & 13,044 & +38 \\
\bottomrule
\end{tabular}
\end{table}

\begin{figure}[H]
    \centering
    \includegraphics[width=0.85\linewidth]{results_a.jpg}
    \caption{Terminal output of the comprehensive model comparison (ReducedMNIST).}
    \label{fig:a_results_terminal}
\end{figure}

\subsubsection{Key Observations}

\begin{itemize}[leftmargin=1.5em]
    \item \textbf{Near-identical accuracy:} Both models achieve $\approx$96.6\%
          test accuracy. The difference ($-$0.07 pp) is negligible and likely within
          run-to-run variance.
    \item \textbf{Winner in test accuracy: Baseline} (96.65\% vs.\ 96.58\%).
    \item \textbf{Winner in training time: Baseline} (11.75\,s vs.\ 21.65\,s) —
          the attention module roughly doubles per-epoch cost on this small model.
    \item \textbf{Generalisation gap:} Both models generalise well (gap $<1\%$),
          indicating that ReducedMNIST is straightforward enough that attention
          provides no regularisation benefit.
    \item \textbf{Minimal parameter overhead:} Only +38 weights ($\approx0.29\%$),
          yet the time cost is disproportionately high due to the extra conv operations
          on every forward pass.
\end{itemize}

\subsection{Analysis}

\subsubsection{Why Attention Makes Little Difference Here}

MNIST digits are centred, well-segmented, and relatively noise-free. The background
pixels carry almost no useful signal, and a simple CNN already learns to ignore them
through weight sharing and pooling. Spatial attention, which is designed to
\emph{suppress} uninformative regions, therefore finds little to improve upon —
the baseline already focuses on the right pixels implicitly.

\subsubsection{Convergence Behaviour}

Both models converge smoothly within the first 2--3 epochs
(Figures~\ref{fig:a_curves_baseline} and~\ref{fig:a_curves_attention}), with training
and validation curves tracking each other closely throughout all 10 epochs. There is
no sign of overfitting, which confirms that the dataset is not challenging enough to
reveal a difference between the two architectures.

\subsubsection{Training Time Overhead}

The attention model takes $\approx$21.65\,s vs.\ 11.75\,s for the baseline — nearly
double. This is because the \texttt{SpatialAttention} module performs two pooling
operations and one additional convolution on every batch, which is relatively costly
when the base model is already very lightweight (only 13,006 parameters).

\subsection{Insights and Observations}

\begin{enumerate}[leftmargin=1.5em]
    \item \textbf{Attention is not universally beneficial.} On clean, centred datasets
          like ReducedMNIST, the baseline already extracts sufficient features, and
          attention adds complexity without improving accuracy.
    \item \textbf{Dataset difficulty matters.} Attention mechanisms tend to shine on
          harder tasks where informative regions are small, sparse, or mixed with
          strong background noise.
    \item \textbf{Parameter efficiency.} The 38-parameter overhead is tiny, but the
          time cost is not — this illustrates that parameter count alone is not a
          reliable proxy for computational cost.
    \item \textbf{Both models generalise well.} Train-test gaps of $<1\%$ indicate
          no meaningful overfitting on ReducedMNIST.
\end{enumerate}

\subsection{Suggestions for Future Improvements}

\begin{itemize}[leftmargin=1.5em]
    \item \textbf{Test on harder datasets} (e.g.\ Fashion-MNIST or CIFAR-10) where
          attention mechanisms are more likely to provide measurable gains.
    \item \textbf{Channel attention (SE blocks):} On MNIST, selecting \emph{which
          feature maps} to emphasise may be more useful than selecting spatial regions.
    \item \textbf{More training epochs:} Both accuracy curves are still slightly
          rising at epoch 10; extending to 20--30 epochs with early stopping may
          reveal a difference.
    \item \textbf{Lighter attention modules:} A $1\times1$ conv instead of $3\times3$
          would reduce the time overhead while preserving the attention mechanism.
\end{itemize}

\newpage
% ══════════════════════════════════════════════════════════════════════════════
%  PART (b)
% ══════════════════════════════════════════════════════════════════════════════
\section{Part (b): CNN vs.\ Attention CNN on Spoken Digits (FSDD)}

\subsection{Introduction}

This part revisits the spoken-digit classification task from Assignment~2. Two CNN
architectures are trained on mel-spectrogram images of spoken digits from the Free
Spoken Digit Dataset (FSDD):

\begin{enumerate}[leftmargin=1.5em]
    \item \textbf{Baseline CNN} — a two-block convolutional classifier with no
          attention mechanism.
    \item \textbf{Attention CNN} — the same architecture augmented with a
          \texttt{SpatialAttentionWithMask} layer after each convolutional block.
\end{enumerate}

\subsection{Dataset and Preprocessing}

\subsubsection{Free Spoken Digit Dataset (FSDD)}

The FSDD contains 3,000 audio recordings (300 per digit, 10 digits), recorded at
8\,kHz by multiple speakers.

\subsubsection{Spectrogram Generation}

\begin{enumerate}[leftmargin=1.5em]
    \item Audio loaded at $f_s = 8{,}000$\,Hz using \texttt{librosa}.
    \item Mel-spectrogram: $N_\text{FFT}=512$, hop $= 256$, $n_\text{mels} = 64$.
    \item Converted to dB: $S_\text{dB} = 10\log_{10}(S/S_\text{max})$.
    \item Min-max normalised and resized to $64\times64$ pixels.
    \item Channel dimension appended: final shape $(64,\,64,\,1)$.
\end{enumerate}

\subsubsection{Train/Test Split}

A stratified 80/20 split (\texttt{random\_state = 42}) gives approximately
2,400 training and 600 test samples, balanced across all ten digit classes.

\subsection{Network Architectures}

\subsubsection{Baseline CNN}

\begin{table}[H]
\centering
\caption{Baseline CNN architecture for FSDD}
\label{tab:b_baseline_arch}
\begin{tabular}{L{5cm} C{3.5cm} C{2.5cm}}
\toprule
\rowcolor{best}
\textbf{Layer (type)} & \textbf{Output Shape} & \textbf{Param \#} \\
\midrule
Conv2D 3$\times$3, 32 filters, ReLU & (None, 64, 64, 32) & 320 \\
\rowcolor{high}
BatchNormalization                   & (None, 64, 64, 32) & 128 \\
MaxPooling2D 2$\times$2              & (None, 32, 32, 32) & 0 \\
\rowcolor{high}
Dropout (0.25)                       & (None, 32, 32, 32) & 0 \\
Conv2D 3$\times$3, 64 filters, ReLU & (None, 32, 32, 64) & 18,496 \\
\rowcolor{high}
BatchNormalization                   & (None, 32, 32, 64) & 256 \\
MaxPooling2D 2$\times$2              & (None, 16, 16, 64) & 0 \\
\rowcolor{high}
Dropout (0.30)                       & (None, 16, 16, 64) & 0 \\
Flatten                              & (None, 16,384)      & 0 \\
\rowcolor{high}
Dense (256, ReLU)                    & (None, 256)         & 4,194,560 \\
Dense (10, Softmax)                  & (None, 10)          & 2,570 \\
\midrule
\rowcolor{best}
\multicolumn{2}{l}{\textbf{Total trainable parameters}} & \textbf{4,216,330} \\
\bottomrule
\end{tabular}
\end{table}

\subsubsection{Spatial Attention Mechanism}

The \texttt{SpatialAttentionWithMask} layer computes:

\begin{equation}
    \mathbf{M} = \sigma\!\left(
        \text{Conv}_{7\times 7}\!\left(
            \left[\,\text{AvgPool}_{C}(\mathbf{F});\;
                  \text{MaxPool}_{C}(\mathbf{F})\,\right]
        \right)
    \right), \qquad
    \mathbf{F}' = \mathbf{F} \odot \mathbf{M}
\end{equation}

Each module uses a $7\times7$ Conv2D with 2 input channels and 1 output channel:
$(7\times7\times2)+1 = \mathbf{99}$ parameters.

\subsubsection{Attention CNN}

\begin{table}[H]
\centering
\caption{Attention CNN architecture for FSDD}
\label{tab:b_attention_arch}
\begin{tabular}{L{5.5cm} C{3.8cm} C{2cm}}
\toprule
\rowcolor{best}
\textbf{Layer (type)} & \textbf{Output Shape} & \textbf{Param \#} \\
\midrule
InputLayer                                       & (None, 64, 64, 1)  & 0 \\
\rowcolor{high}
Conv2D 3$\times$3, 32 filters, ReLU             & (None, 64, 64, 32) & 320 \\
BatchNormalization                               & (None, 64, 64, 32) & 128 \\
\rowcolor{mid}
\textbf{SpatialAttentionWithMask (attn\_1)}      & [(None,64,64,32), (None,64,64,1)] & \textbf{99} \\
MaxPooling2D 2$\times$2                          & (None, 32, 32, 32) & 0 \\
\rowcolor{high}
Conv2D 3$\times$3, 64 filters, ReLU             & (None, 32, 32, 64) & 18,496 \\
BatchNormalization                               & (None, 32, 32, 64) & 256 \\
\rowcolor{mid}
\textbf{SpatialAttentionWithMask (attn\_2)}      & [(None,32,32,64), (None,32,32,1)] & \textbf{99} \\
MaxPooling2D 2$\times$2                          & (None, 16, 16, 64) & 0 \\
\rowcolor{high}
Flatten                                          & (None, 16,384)      & 0 \\
Dense (256, ReLU)                                & (None, 256)         & 4,194,560 \\
\rowcolor{high}
Dense (10, Softmax)                              & (None, 10)          & 2,570 \\
\midrule
\rowcolor{best}
\multicolumn{2}{l}{\textbf{Total trainable parameters}} & \textbf{4,216,528} \\
\bottomrule
\end{tabular}
\end{table}

\subsubsection{Parameter Overhead of Attention}

Two attention modules add $2\times99=+198$ parameters over the baseline, an
increase of only $\approx 0.005\%$.

\subsection{Training Configuration}

\begin{table}[H]
\centering
\caption{Training hyperparameters for Part (b)}
\label{tab:b_hyperparams}
\begin{tabular}{L{5.5cm} C{5cm}}
\toprule
\rowcolor{best}
\textbf{Hyperparameter} & \textbf{Value} \\
\midrule
Optimizer                       & Adam (lr = 0.001) \\
\rowcolor{high}
Loss function                   & Sparse categorical cross-entropy \\
Batch size                      & 32 \\
\rowcolor{high}
Epochs                          & 10 \\
Train / Test split              & 80\% / 20\% (stratified) \\
\rowcolor{high}
Input image size                & 64 $\times$ 64 $\times$ 1 \\
Number of classes               & 10 (digits 0--9) \\
\rowcolor{high}
Sample rate                     & 8,000 Hz \\
FFT window ($N_\text{FFT}$)     & 512 \\
\rowcolor{high}
Hop length                      & 256 \\
Mel bins                        & 64 \\
\bottomrule
\end{tabular}
\end{table}

\subsubsection{Training Logs}

\begin{figure}[H]
    \centering
    \includegraphics[width=\linewidth]{training_baseline.jpg}
    \caption{Baseline CNN training log (FSDD).}
    \label{fig:b_log_baseline}
\end{figure}

\begin{figure}[H]
    \centering
    \includegraphics[width=\linewidth]{training_attention.jpg}
    \caption{Attention CNN training log (FSDD).}
    \label{fig:b_log_attention}
\end{figure}

\subsection{Results}

\subsubsection{Training and Validation Curves}

\begin{figure}[H]
    \centering
    \includegraphics[width=\linewidth]{Figure_5.png}
    \caption{Baseline CNN accuracy and loss curves (FSDD).}
    \label{fig:b_curves_baseline}
\end{figure}

\begin{figure}[H]
    \centering
    \includegraphics[width=\linewidth]{Figure_6.png}
    \caption{Attention CNN accuracy and loss curves (FSDD).}
    \label{fig:b_curves_attention}
\end{figure}

\subsubsection{Quantitative Comparison}

\begin{table}[H]
\centering
\caption{Final-epoch performance: Baseline vs.\ Attention CNN (FSDD)}
\label{tab:b_results}
\begin{tabular}{L{4.2cm} C{2.8cm} C{2.8cm} C{2.5cm}}
\toprule
\rowcolor{best}
\textbf{Metric} & \textbf{Baseline CNN} & \textbf{Attention CNN} & \textbf{Difference} \\
\midrule
Train Accuracy (\%)   & 97.13 & 99.21 & \textcolor{softblue}{+2.08} \\
\rowcolor{high}
Test Accuracy (\%)    & 93.67 & 98.17 & \textcolor{softblue}{+4.50} \\
Train Loss            & 0.1247 & 0.0260 & \textcolor{softblue}{$-$0.0988} \\
\rowcolor{high}
Test Loss             & 0.3967 & 0.0667 & \textcolor{softblue}{$-$0.3300} \\
Generalisation Gap    & 3.46\% & 1.04\% & \textcolor{softblue}{$-$2.42 pp} \\
\rowcolor{high}
Model Parameters      & 4,216,330 & 4,216,528 & +198 \\
\bottomrule
\end{tabular}
\end{table}

\begin{figure}[H]
    \centering
    \includegraphics[width=0.85\linewidth]{results.jpg}
    \caption{Terminal output of the comprehensive model comparison (FSDD).}
    \label{fig:b_results_terminal}
\end{figure}

\subsubsection{Key Observations}

\begin{itemize}[leftmargin=1.5em]
    \item \textbf{Better test accuracy:} Attention CNN achieves \textbf{98.17\%}
          vs.\ 93.67\% — a gain of \textbf{4.50 pp}.
    \item \textbf{Lower test loss:} Attention model test loss (0.0667) is nearly
          \textbf{6$\times$} smaller than the baseline (0.3967).
    \item \textbf{Better generalisation:} Train-test gap drops from 3.46\% to 1.04\%,
          showing spatial attention acts as an implicit regulariser.
    \item \textbf{Negligible overhead:} Only +198 parameters ($\approx 0.005\%$).
\end{itemize}

\subsection{Analysis}

\subsubsection{Why Attention Helps on Spectrograms}

Mel-spectrograms have a clear spatial structure — frequency on the vertical axis,
time on the horizontal. Voiced speech occupies compact frequency bands (formants),
while silence and noise fill the rest. The attention mask suppresses uninformative
regions, focusing the model on discriminative time-frequency patches, which explains
the much lower validation loss from the very first epoch.

\subsubsection{Convergence Behaviour}

\begin{itemize}[leftmargin=1.5em]
    \item The \textbf{baseline} validation loss spikes dramatically in epochs 1--3
          ($\approx$63.5), indicating memorisation before generalisation.
    \item The \textbf{attention} model's peak loss reaches only $\approx$10.7,
          converging far more smoothly to a superior minimum.
\end{itemize}

\subsubsection{Generalisation Gap}

The attention model's smaller train-test gap (1.04\% vs.\ 3.46\%) confirms reduced
overfitting. The spatial mask performs feature selection in the spatial domain,
acting analogously to Dropout but operating on feature-map positions rather than
individual neurons.

\subsubsection{Training Time}

Per-step time is marginally higher for the attention model ($\sim$135--146\,ms/step
vs.\ $\sim$106--113\,ms/step for the baseline), amounting to roughly 2--4 seconds
per epoch — negligible given the accuracy gains.

\subsection{Insights and Observations}

\begin{enumerate}[leftmargin=1.5em]
    \item \textbf{Spatial attention suits spectrogram inputs.} The grid-like structure
          of mel-spectrograms with localised time-frequency information makes them
          ideal for spatial attention.
    \item \textbf{Two attention layers outperform one.} Applying attention after both
          the 32-filter and 64-filter blocks allows refocusing at different abstraction
          levels.
    \item \textbf{BatchNorm complements attention.} Normalising feature distributions
          while reweighting spatial positions yields both stable training and selective
          focus.
    \item \textbf{The Dense bottleneck dominates parameter count.} Over 99\% of
          parameters reside in the first Dense layer; global average pooling would
          remove this bottleneck.
    \item \textbf{Variable-length audio mapped to fixed grids.} All recordings are
          compressed to 64$\times$64; attention partially compensates by focusing on
          the most informative temporal segments.
\end{enumerate}

\subsection{Suggestions for Future Improvements}

\begin{itemize}[leftmargin=1.5em]
    \item \textbf{CBAM:} Combines channel attention with spatial attention for
          stronger feature selection.
    \item \textbf{Self-attention / Transformer blocks:} Captures long-range
          dependencies, useful for digits with similar formant patterns.
    \item \textbf{Squeeze-and-Excitation (SE) networks:} Channel-wise recalibration
          with minimal parameter overhead.
    \item \textbf{Global Average Pooling (GAP):} Replace Flatten $\to$ Dense(256)
          with GAP $\to$ Dense(10) to reduce parameters from 4.2M to $\sim$650.
    \item \textbf{SpecAugment:} Masking random frequency bands and time steps to
          improve robustness to speaker variability.
    \item \textbf{More epochs with early stopping:} The attention accuracy curve
          is still improving at epoch 10; 20--30 epochs with patience $= 5$ could
          push test accuracy above 99\%.
    \item \textbf{Speaker-independent cross-validation:} Leave-one-speaker-out
          evaluation for a more realistic generalisation estimate.
\end{itemize}

% ══════════════════════════════════════════════════════════════════════════════
%  CROSS-PART COMPARISON
% ══════════════════════════════════════════════════════════════════════════════
\section{Cross-Part Comparison: When Does Attention Help?}

\begin{table}[H]
\centering
\caption{Summary comparison of attention impact across both parts}
\label{tab:cross_comparison}
\begin{tabular}{L{4cm} C{3cm} C{3cm} C{3cm}}
\toprule
\rowcolor{best}
\textbf{Aspect} & \textbf{Part (a) MNIST} & \textbf{Part (b) FSDD} & \textbf{Winner} \\
\midrule
Dataset type            & Clean images   & Spectrograms   & -- \\
\rowcolor{high}
Attention benefit       & Negligible     & Significant    & Part (b) \\
Test accuracy gain      & $-$0.07 pp     & +4.50 pp       & Part (b) \\
\rowcolor{high}
Generalisation improvement & Slight decrease & Large decrease & Part (b) \\
Parameter overhead      & +38 (0.29\%)   & +198 (0.005\%) & Part (b) \\
\rowcolor{high}
Time overhead           & $+$9.9 s ($\times$1.8) & Minor ($\sim$2--4 s/ep) & Part (b) \\
\bottomrule
\end{tabular}
\end{table}

The comparison reveals a clear pattern: \textbf{spatial attention provides the most
benefit when the input contains localised, structured signal embedded in noise} —
exactly the case for mel-spectrograms of spoken digits. On clean, centred images
(ReducedMNIST), the baseline CNN is already sufficient and attention adds only
time cost without accuracy gains.

\end{document}
